\chapter{Diagonalization}
Given a linear operator $T$ on a finite dimensional vector space $V$, the \emph{diagonalization problem} asks whether an ordered basis $\beta$ exists such that $[T]_\beta$ is diagonal, and how to find it if so.
Since diagonal matrices are computationally simple, this helps us understand better how $T$ acts on $V$ and solve many practical problems that can be posed with linear algebra.
Eigenvalues and eigenvectors are essential to the diagonalization problem and in their own right.
In the study of mechanical systems, one must control unwanted vibrations to prevent failures, eg a bridge collapse. These involve eigenvalues and eigenvectors of the differential equations that model the system.
Other applications are facial recognition, fingerprint authentication, and principal components analysis.

\section{Eigenvalues and eigenvectors}
Recall in Example \ref{e:reflect} that reflection across the line $y=2x$ is represented by a diagonal matrix with respect to the basis $\{ (1,2)^\top,(-2,1)^\top \}$.
\begin{defn}
A linear operator $T$ on a finite dimensional vector space $V$ is \emph{diagonalizable} if an ordered basis $\beta$ exists such that $[T]_\beta$ is diagonal. $A \in F^{n \times n}$ is diagonalizable if $L_A$ is.\\
If $V$ is a (not necessarily finite dimensional) vector space, a nonzero $\mb{v}\in V$ is an \emph{eigenvector} of $T$ with \emph{eigenvalue} $\lambda$ if $T( \mb{v})= \lambda \mb{v}$. $A \in F^{n \times n}$ has (nonzero) eigenvector $\mb{v}\in F^n$ with eigenvalue $\lambda$ if $L_A$ does.\\
Eigenvectors (resp values) are sometimes called characteristic or proper vectors (resp values).
\end{defn}
\begin{lem}\label{l:diagze-basis}
A linear operator $T$ on a finite dimensional vector space $V$ is diagonalizable if and only if $V$ has an ordered basis $\beta = \{ \mb{v}_1, \ldots, \mb{v}_n \}$ consisting of eigenvectors of $T$. If so, $[T]_\beta$ is diagonal with the corresponding eigenvalues on its main diagonal.\\
$A \in F^{n \times n}$ is diagonalizable if and only if it has an ordered basis $\beta = \{ \mb{v}_1, \ldots, \mb{v}_n \}$ consisting of eigenvectors of $A$. If these eigenvectors form the columns of $Q$, then $Q^{-1}AQ$ is diagonal with the corresponding eigenvalues on its main diagonal. $A$ is thus diagonalizable if and only if it is similar to a diagonal matrix.\\
To \emph{diagonalize} $A$ is to find an eigenbasis and the corresponding eigenvalues.
\end{lem}
Here is our main tool for finding eigenvalues.
\begin{lem}\label{l:chr-poly}
$A \in F^{n \times n}$ has eigenvalue $\lambda \in F$ if and only if $\det (A- \lambda I_n)=0$. $t \mapsto \det (A-tI_n)$ is the \emph{characteristic polynomial} of $A$.\footnote{technically each entry of $A-tI_n$ is not a scalar, but a member of $F[t]$, the field of quotients of polynomials in $t$ with coefficients in $F$}
If $A$ and $B$ are similar matrices, then $\det A= \det B$ (and hence have the same characteristic polynomial) and $\tr A= \tr B$.
\end{lem}
\begin{proof}%easy, short
$A$ has eigenvalue $\lambda$ iff there exists nonzero $\mb{v}\in F^n$ with $A \mb{v}= \lambda \mb{v}$ (this $\mb{v}$ is an eigenvector) iff $\rank (A- \lambda I_n)<n$ (Lemma \ref{l:hom-ker}) iff $\det (A- \lambda I_n)=0$ (Lemma \ref{l:det-mult} part 2 and Lemma \ref{l:mat-inv-1to1}).\\
Pick invertible $Q$ such that $B=Q^{-1}AQ$, then  $B- tI_n=Q^{-1}(A-tI_n)Q$ also, and $\tr B= \tr AQQ^{-1}= \tr A$ (Exercise \ref{e:tr-commute}), $\det B= \det Q \det A \det Q^{-1}= \det A$ (Lemma \ref{l:det-mult}), and likewise $\det(B- tI_n)= \det(A-tI_n)$.
\end{proof}
We can use the fact that similar matrices have the same determinant and characteristic polynomial to define those things for linear operators.
\begin{defn}\label{d:op-det}
Let $V$ be a vector space with $\dim V=n \in \bb{Z}^+$, $T$ a linear operator on $V$, $\beta$ any ordered basis of $V$. Then define $\det T:= \det [T]_\beta$ and the characteristic polynomial of $T$ to be $t \mapsto \det ([T]_\beta -I_n)$.
\end{defn}
These are well defined: if $\gamma$ is another ordered basis of $V$, then $[T]_\beta$ and $[T]_\gamma$ are similar (Lemma \ref{l:change-coord}).
\begin{lem}\label{l:chrpoly-degn}
For $A \in F^{n \times n}$, its characteristic polynomial has degree $n$ and lead coefficient $(-1)^n$. In fact, it has the form $\prod_{j=1}^n(a_{jj}-t)+q(t)$ where $q$ has degree $n-2$.
$A$ thus has at most $n$ distinct eigenvalues. We also have the coefficient of $t^{n-1}$ as $(-1)^{n-1}\tr A$.
$A$ and $A^\top$ have the same characteristic polynomial (and hence the same eigenvalues).
\end{lem}
\begin{proof}
Put $\tilde{I}_n:=I_n- \mb{e}_1 \mb{e}_1^\top$ (the identity matrix, except with its 1,1 entry 0 instead of 1). We use induction on $n$.
The base case $n=2$ is clear by direct computation: $\det (A-tI_2)=(a_{11}-t)(a_{22}-t)-a_{12}a_{21}$.
For $n>2$ we have $\det (A-tI_n)=(a_{11}-t) \det \overline{(A-tI_n)}_{11}+ \sum_{k=2}^n(-1)^{k+1}a_{1k}\det \overline{(A-tI_n)}_{1k}$.
The induction hypothesis applies directly to $\det \overline{(A-tI_n)}_{11}= \prod_{j=2}^n(a_{jj}-t)+ \tilde{q}(t)$ (where $\tilde{q}$ has degree $n-3$). And for $k>1$, $\overline{(A-tI_n)}_{1k}$ has the form $B-t \tilde{I}_{n-1}$ after rearranging rows which affects only the sign by Lemma \ref{l:det-elem} part 6.
A parallel induction argument shows $\det (B-t \tilde{I}_{n-1})$ has degree $n-2$. Putting these together gives the claimed form that $\det (A-tI_n)$ takes.\\
By Corollary \ref{c:polyn0}, this polynomial has at most $n$ distinct roots. By Lemma \ref{l:chr-poly}, $A$ has at most $n$ distinct eigenvalues.\\
Finally, we have $\det (A^\top -tI_n)= \det (A-tI_n)^\top = \det (A-tI_n)$ by Lemma \ref{l:det-mult} part 3.
\end{proof}
\begin{example}
Not all operators/ matrices are diagonalizable. By geometric intuition, the rotation operator by $\theta$ on $\bb{R}^2$ has no eigenvalues in $\bb{R}$ except when the rotation is by $0^\circ$ or $180^\circ$.
The characteristic polynomial is $t^2-2t \cos \theta +1$, which has no roots in $\bb{R}$ when $0< \theta < \pi$.\\
By Lemma \ref{l:diff-eq-dim1}, the differential operator $D$ on $C^\infty (\bb{R}, \bb{C})$ has any $\lambda \in \bb{C}$ as an eigenvalue, with eigenvector $t \mapsto e^{\lambda t}$.
\end{example}
We reduce the problem of finding eigenvectors of a linear operator tp that for a matrix.
\begin{exer}
Let $V$ be a vector space with $\dim V=n \in \bb{Z}^+$, $T$ a linear operator on $V$, $\beta$ any ordered basis of $V$.
Then $\mb{v}\in V$ is an eigenvector of $T$ with eigenvalue $\lambda$ if and only if $[ \mb{v}]_\beta$ is an eigenvector of $[T]_\beta$ with eigenvalue $\lambda$.
An equivalent restatemeht using $\phi_\beta:V \to F^n$, $\mb{v}\mapsto [ \mb{v}]_\beta$ is: $\mb{y}\in F^n$ is an eigenvector of $[T]_\beta$ with eigenvalue $\lambda$ if and only if $\phi_\beta^{-1}( \mb{y})$ is an eigenvector of $T$ with eigenvalue $\lambda$.
\end{exer}
\begin{proof}%easy, short
By Lemma \ref{l:std-repr-iso}, $\phi_\beta$ is an isomorphism. We have $T( \mb{v})= \lambda \mb{v}$ $\leftrightarrow$ $[T( \mb{v})]_\beta =[ \lambda \mb{v}]_\beta$ $\leftrightarrow$ $[T]_\beta [ \mb{v}]_\beta = \lambda[ \mb{v}]_\beta$, where LHS uses Proposition \ref{p:transform-coordinate-vector} and RHS uses $\phi_\beta$ is an isomorphism (in particular, linear).
We can replace $[ \mb{v}]_\beta$ with $\mb{y}$ and $\mb{v}$ with $\phi_\beta^{-1}( \mb{y})$.
\end{proof}
\begin{exer}\label{e:eig-poly}
Let $T$ be a linear operator on a vector space $V$.
\begin{itemize}
\item $T$ is invertible if and only if 0 is not an eigenvalue, in which case $\mb{v}\in V$ is an eigenvector of $T$ with eigenvalue $\lambda$ if and only if $\mb{v}$ is an eigenvector of $T^{-1}$ with eigenvalue $1/ \lambda$.
\item If $p \in P(F)$, and $\mb{v}\in V$ is an eigenvector of $T$ with eigenvalue $\lambda$, then $\mb{v}$ is an eigenvector of $p(T)$ with eigenvalue $p( \lambda)$.
\item If $V$ is finite dimensional and $f$ is the characteristic polynomial of $T$, then $f(T)$ is the zero operator. \footnote{this is a special case of a powerful theorem we see later}
\end{itemize}
Analogous results for all of the above hold for matrices.
\end{exer}
\begin{proof}% easy, short
Part 1: By definition, 0 is not an eigenvalue of $T$ $\leftrightarrow$ $\ker T= \{ \mb{0}_V \}$ $\leftrightarrow$ $T$ is invertible by Lemma \ref{l:inv-1to1-or-onto}.
We have $\mb{v}=T^{-1}\circ T( \mb{v})= \lambda T^{-1}( \mb{v})$ by Lemma \ref{l:inv-linear} $leftrightarrow$ $T^{-1}( \mb{v})= \mb{v}/ \lambda$.\\
Part 2: We 1st show the special case $T^m( \mb{v})= \lambda^m \mb{v}$ for $m \in \bb{Z}_{\geq 0}$ by induction. The base case $m=0$ is clear since $T^0=I_V$.
When $m>0$, $T^m( \mb{v})=T \circ T^{m-1}( \mb{v})=T( \lambda^{m-1}\mb{v})= \lambda^m \mb{v}$ by the inductive hypothesis then linearity.
For general $p \in P(F)$, say $p(t)= \sum_{j=0}^ma_jt^j$. Then $p(T)( \mb{v})= \sum_{j=0}^ma_jT^j( \mb{v})= \sum_{j=0}^ma_j \lambda^j \mb{v}=p( \lambda) \mb{v}$.\\
Part 3: By Lemma \ref{l:diagze-basis}, we can pick a basis $\beta$ of $V$ of eigenvectors of $T$. Say $\mb{v}\in \beta$ with eigenvalue $\lambda$.
Then by part 2, $f(T)( \mb{v})=f( \lambda)( \mb{v})= \mb{0}_V$ by Lemma \ref{l:chr-poly}. $f(T)$ is thus the zero operator by Corollary \ref{c:lin-map-basis}.
\end{proof}
When $V= \bb{R}^n$ and $F= \bb{R}$, and $T$ is a linear operator on $\bb{R}^n$ with eigenvector $\mb{v}$ and eigenvalue $\lambda$, we can describe geometrically how $T$ acts on $\Span \{  \mb{v}\}$.
By linearity, for any $c \in \bb{R}$, $T(c \mb{v})=c \lambda \mb{v}$. Fix $\mb{w}\in \Span \{  \mb{v}\}$. When $\lambda>1$, $\mb{w}$ is moved farther from the origin by a factor of $\lambda$.
When $\lambda =1$, $\mb{w}$ is unchanged ($T$ acts like the identity operator). When $0< \lambda <1$, $\mb{w}$ is moved closer to the origin by a factor of $\lambda$.
When $\lambda =0$, $T$ acts like the zero operator. When $\lambda <0$, $T$ flips $\mb{w}$ about the origin and scales by $| \lambda |$ as previously described.

\section{Diagonalizability}
Though Lemma \ref{l:diagze-basis} gives a necessary and sufficient condition for diagonalizability, and Lemma \ref{l:chr-poly} helps find eigenvalues and their associated eigenvalues, we still need a test for diagonalizability, and a way to find an eigenvector basis in order to solve the diagonalization problem.
While finding eigenvectors for each eigenvalue is not enough to result in a basis, they are at least LID.
\begin{lem}\label{l:dnct-eval-lid}
Let $T$ be a linear operator on vector space $V$. If $\lambda_1, \ldots, \lambda_k$ are distinct eigenvalues of $T$ and $S_r$ is a LID set of eigenvectors corresponding to $\lambda_r$ for $r=\ot{k}$, then $\cup_{r=1}^kS_r$ is LID.
In particular, if $\dim V=n$ and $T$ has $n$ distinct eigenvalues, then $T$ is diagonalizable.
\end{lem}
\begin{proof}%rating 1/2
Use induction on $k$. For the base case $k=1$, $S_1$ is LID by construction. For $k>1$, we assume the linear independence of $\cup_{r=1}^{k-1}S_r= \{ \mb{v}_1, \ldots, \mb{v}_n \}$ with corresponding eigenvalues $\lambda(j)$ for $j= \ot{n}$, none of which are $\lambda_k$.
Say $S_k= \{ \mb{w}_1, \ldots, \mb{w}_m \}$, and suppose we have coefficients $a_j,b_j \in F$ such that $\sum_{j=1}^na_j \mb{v}_j+ \sum_{j=1}^mb_j \mb{w}_j= \mb{0}$, that is $\sum_{j=1}^na_j \mb{v}_j=- \sum_{j=1}^mb_j \mb{w}_j$.
Then $T \big( \sum_{j=1}^mb_j \mb{w}_j \big) = \lambda_k \sum_{j=1}^mb_j \mb{w}_j=- \lambda_k \sum_{j=1}^na_j \mb{v}_j$. But also $T \big( \sum_{j=1}^mb_j \mb{w}_j \big) =T \big( - \sum_{j=1}^na_j \mb{v}_j \big) =- \sum_{j=1}^na_j \lambda(j) \mb{v}_j$.
Putting these 2 pieces together, $\sum_{j=1}^na_j \{ \lambda_k- \lambda(j) \} \mb{v}_j= \mb{0}_V$. Since $\cup_{r=1}^{k-1}S_r$ is LID, we have each $a_j \{ \lambda_k- \lambda(j) \} =0$ and in fact each $a_j=0$ since $\lambda(j) \neq \lambda_k$.
Since $S_k$ is LID, we also have each $b_j=0$, so $\cup_{r=1}^kS_r$ is LID.\\
When $\dim V=n$ and $T$ has $n$ distinct eigenvalues, we can pick 1 eigenvector corresponding to each. This set of $n$ eigenvectors is LID, and hence a basis by Corollary \ref{c:extend-LID2basis}.
\end{proof}
Having $n$ distinct eigenvalues is not a necessary condition for diagonalizability. The identity operator/ matrix is diagonalizable, despite having just 1 distinct eigenvalue.
\begin{defn}
$p \in P(F)$ \emph{splits} over $F$ if there exist $c,a_1, \ldots,a_n \in F$, not necessarily distinct, such that $p(t)=c \prod_{j=1}^n(t-a_j)$.
\end{defn}
$p(t)=t^2+1$ splits over $\bb{C}$ but not $\bb{R}$.
\begin{lem}
If linear operator $T$ on vector space $V$ over field $F$ is diagonalizable, then its characteristic polynomial splits over $F$.
\end{lem}
\begin{proof}%easy, short
By Lemma \ref{l:diagze-basis}, pick a basis $\beta$ for $V$ of eigenvectors of $T$. Then $[T]_\beta = \T{diag}( \lambda_1,..., \lambda_n)$ where each $\lambda_j$ is an eigenvalue of $T$.
Then the characteristic polynomial is $t \mapsto \prod_{j=1}^n( \lambda_j-t)$.
\end{proof}
Thus, if $T$ is diagonalizable but does not have $n$ distinct eigenvalues, then some eigenvalue $\lambda$ shows up multiple times in the characteristic polynomial.
\begin{example}\label{e:splits-not-diagle}
The differential operator on $P_2( \bb{R})$ has characteristic polynomial $t \mapsto -t^3$, which splits over $\bb{R}$. Its only eigenvalue is 0, with the constant functions as eigenvectors, so it is not diagonalizable.
\end{example}
Since an eigenvector with eigenvalue $\lambda$ is a member of $\ker (T- \lambda I_V)$, this subspace is naturally of interest.
\begin{defn}
If $f$ is the characteristic polynomial for an operator $T$ on $V$ (or matrix) with eigenvalue $\lambda$, its \emph{(algebraic) multiplicity} is the largest $k \in \bb{Z}^+$ such that $t \mapsto (t- \lambda)^k$ divides $f$.\\
The \emph{eigenspace} of $T$ for eigenvector $\lambda$ is $E_\lambda := \ker (T- \lambda I_V)$
\end{defn}
\begin{prop}\label{p:diagze-mult}
Let $T$ be a linear operator on a vector space $V$ of dimension $n \in \bb{Z}^+$, with distinct eigenvalues $\lambda_1, \ldots, \lambda_k$.
\begin{itemize}
\item For any eigenvalue $\lambda \in \{ \lambda_1, \ldots, \lambda_k \}$ with multiplicity $m$, we have $1 \leq \dim E_\lambda \leq m$.
\item When the characteristic polynomial splits over $F$, $T$ is diagonalizable if and only if $\dim E_\lambda =m$ for all $\lambda \in \{ \lambda_1, \ldots, \lambda_k \}$, where $m$ is its multiplicity.
$T$ is also diagonalizable if and only if $\cup_{j=1}^k \beta_j$ is an eigenvector basis for $V$, where each $\beta_j$ is a basis for the eigenspace corresponding to $\lambda_j$.
\end{itemize}
\end{prop}
\begin{proof}%easy, short
Part 1: Put $s:= \dim E_\lambda$. Any eigenvector corresponding to $\lambda$ is a member of $E_\lambda$, so $s \geq 1$. Pick a basis $\{ \mb{v}_1, \ldots, \mb{v}_s \}$ of $E_\lambda$, extend it to a basis $\beta = \{ \mb{v}_1, \ldots, \mb{v}_n \}$ of $V$ (Corollary \ref{c:extend-LID2basis}).
Then $[T]_\beta$ has the form as in Exercise \ref{e:block-det}. The characteristic polynomial is
\[ \det ([T]_\beta -tI_n)= \det \tbtm{( \lambda -t)I_s}{O}{*}{*}\]
So $t \mapsto (\lambda -t)^s$ divides the characteristic polynomial, so $s \leq m$.\\
Part 2: "$\to$" Use Lemma \ref{l:diagze-basis} to pick a basis $\beta = \{ \mb{v}_1, \ldots, \mb{v}_n \}$ of eigenvectors of $T$.
Without loss of generality, the 1st $m$ correspond to eigenvalue $\lambda$, and are LID in $E_\lambda$, so $s \geq m$ (Corollary \ref{c:extend-LID2basis}).
"$\leftarrow$ Since the characteristic polynomial splits and it has degree $n$ (Lemma \ref{l:chrpoly-degn}), the sum of all multiplicities over $\lambda_1, \ldots, \lambda_k$ is $n$.
By hypothesis, this is also the sum of the eigenspaces' dimensions, so $\cup_{j=1}^k \beta_j$ has $n$ LID eigenvectors (Lemma \ref{l:dnct-eval-lid}), hence is a basis of $V$ (Corollary \ref{c:extend-LID2basis}). This also shows $T$ is diagonalizable (Lemma \ref{l:diagze-basis}).
\end{proof}
Proposition \ref{p:diagze-mult} gives the way to solve the diagonalization problem. To summarize, find the characteristic polynomial (sometimes there is a convenient 'starting' basis). If it splits, check for each repeated eigenvalue whether $\dim E_\lambda$ is its multiplicity $m$ (this is automatically true when $m=1$).
If so, then $T$ is diagonalizable. Find bases for each eigenspace and take the union to get an eigenvector basis for $T$.
It is easy to take powers of diagonalizable operators/ matrices since to find $D^m$ for diagonal $D \in F^{n \times n}$ we just take powers along the main diagonal. $[T^m]_\beta=([T]_\beta)^m$ (Lemma \ref{l:matrix-compose} and induction) is easy to compute when $[T]_\beta$ is diagonal.
By Lemma \ref{l:diagze-basis}, if $A \in F^{n \times n}$ is diagonalizable, then we can write $D=Q^{-1}AQ$ $\leftrightarrow$ $A=QDQ^{-1}$ where $Q$ has eigenvectors of $A$ as its columns and $D$ is diagonal with the corresponding eigenvalues along its main diagonal. Induction shows $A^m=QD^mQ^{-1}$.
\begin{prop}
If $A \in \bb{C}^{n \times n}$ is diagonalizable and $\mb{y}: \bb{R}\to \bb{C}^n$ is differentiable, then the general solution to the system of differential equations $\mb{y}'=A \mb{y}$ has the form $\mb{y}(t)= \sum_{j=1}^n \mb{u}_j \exp( \lambda_jt)$, where each $\mb{u}_j$ can be any vector from the $\lambda_j$ eigenspace.\footnote{so $\mb{u}_j= \mb{0}_n$ is allowed}
When $A \in \bb{R}^{n \times n}$, if $\lambda = \alpha + \beta i$ is an eigenvalue with eigenvector $\mb{u}= \mb{a}+ \mb{b}i$, then $\bar{\lambda}$ is also an eigenvalue with eigenvector $\bar{\mb{u}}$, and we can replace the 2 terms corresponding to them in the general form with $c_1e^{\alpha t}( \mb{a}\cos \beta t- \mb{b}\sin \beta t)+c_2e^{\alpha t}( \mb{a}\sin \beta t+ \mb{b}\cos \beta t)$ for any $c_1, c_2 \in \bb{R}$.
\end{prop}
\begin{proof}%see also Nagle, sections 9.5 and 9.6
Use Lemma \ref{l:diagze-basis} to write $D=Q^{-1}AQ$ where $Q$ has eigenvectors $\mb{u}_j$ of $A$ as its columns and $D$ is diagonal with the corresponding eigenvalues $\lambda_j$ along its main diagonal.\footnote{the only difference from the general solution is that here $\mb{u}_j \neq \mb{0}_n$}
Left multiply both sides by $Q^{-1}$ and use the linearity of the derivative to get $\mb{y}'=A \mb{y}$ $\leftrightarrow$ $(Q^{-1}\mb{y})'=Q^{-1}\mb{y}'=DQ^{-1} \mb{y}$.
Put $\mb{s}:=Q^{-1}\mb{y}$. Then we have $s'_j= \lambda s_j$ for $j=\ot{n}$. All solutions have the form $s_j(t)=c_j \exp( \lambda_jt)$ for $c_j \in \bb{C}$ (Lemma \ref{l:diff-eq-dim1}).
Using  $\mb{y}=Q \mb{s}$, all solutions have the form $\mb{y}(t)= \sum_{j=1}^nc_j \mb{u}_j \exp( \lambda_jt)$, and each $c_j \mb{u}_j$ is in the $\lambda_j$ eigenspace.
The functions $\mb{u}_j \exp( \lambda_jt)$ span the solution space. They are also LID since the $\mb{u}_j$ are and $\exp( \lambda_jt) \neq 0$.\\
When $A \in \bb{R}^{n \times n}$, 1st note that taking complex conjugates on both sides of $A \mb{u}= \lambda \mb{u}$ gives $A \bar{\mb{u}}= \bar{\lambda}\bar{\mb{u}}$, so eigenvalues and eigenvectors come in conjugate pairs.
We can replace $\{ \mb{u}e^{ \lambda t}, \bar{\mb{u}}e^{ \bar{\lambda}t}\}$ with $\{ \Re \mb{u}e^{ \lambda t}, \Im \mb{u}e^{ \lambda t}\} = \{ e^{\alpha t}( \mb{a}\cos \beta t- \mb{b}\sin \beta t), e^{\alpha t}( \mb{a}\sin \beta t+ \mb{b}\cos \beta t) \}$.
\end{proof}
\begin{example}
For the system $y'_1=y_1+y_2$, $y'_2=2y_1+2y_2$, the general solution is
\[ \vect{y_1(t)}{y_2(t)}=c_1 \vect{1}{2}e^{3t}+c_2 \vect{1}{-1}\]
so $A$ need not be invertible, in which case there are additive constant terms.\\
For the system $y'_1=y_1-y_2$, $y'_2=y_1+y_2$, the coefficient matrix has eigenvalue $1-i$ with eigenvector $(1,i)^\top$ (and their complex conjugates), so the general solution is
\[ \vect{y_1(t)}{y_2(t)}=e^t \big\{ c_1 \vect{\cos t}{\sin t}+c_2 \vect{\sin t}{- \cos t}\big\} \]
For the system $y'_1=y_1$, $y'_2=y_2$, the coefficient matrix is $I_2$. We can pick any 2 LID vectors to construct the general solution
\[ \vect{y_1(t)}{y_2(t)}=c_1 \vect{1}{0}e^t+c_2 \vect{1}{1}e^t \]
While not a very interesting example, this is just to reassure that it does not matter what basis is used in the case of repeating eigenvalues.
\end{example}
\begin{defn}
If $V$ is a vector space with subspaces $W_1, \ldots,W_k$, we define their \emph{sum} to be $\sum_{j=1}^kW_j:= \{ \sum_{j=1}^k \mb{v}_j: \mb{v}_j \in W_j \}$. It is a \emph{direct sum}, written $\bigoplus_{j=1}^kW_j$, if it is a sum, and $W_j \cap \sum_{s \neq j}W_s= \{ \mb{0}_V \}$ for all $j= \ot{k}$.
\end{defn}\label{l:directsum}
Several characterizations of the direct sum are useful.
\begin{lem}
In the above definition, $V= \bigoplus_{j=1}^kW_j$ is equivalent to any of:
\begin{itemize}
\item $V= \sum_{j=1}^kW_j$, and if $\mb{v}_j \in W_j$ and $\sum_{j=1}^k \mb{v}_j= \mb{0}_V$, then each $\mb{v}_j= \mb{0}_V$.
\item Any $\mb{v} \in V$ can be uniquely written $\mb{v}= \sum_{j=1}^k \mb{v}_j$ where $\mb{v}_j \in W_j$.
\item If $\gamma_j$ is a basis for $W_j$, then $\cup_{j=1}^k \gamma_j$ is a basis for $V$. Equivalently, $\dim V= \sum_{j=1}^k \dim W_j$.
\item For each $j= \ot{k}$ there exists a basis $\gamma_j$ for $W_j$ such that $\cup_{j=1}^k \gamma_j$ is a basis for $V$.
\end{itemize}
\end{lem}
\begin{proof}
Direct sum $\to$ 1: We have $- \mb{v}_1 \in W_1 \cap \sum_{j=2}^kW_j$, so $\mb{v}_j= \mb{0}_V$, same with the other indices.
1 $\to$ 2: if $\mb{v}_j, \mb{w}_j \in W_j$ are such that $\sum_{j=1}^k \mb{v}_j= \sum_{j=1}^k \mb{w}_j$, then $\sum_{j=1}^k( \mb{v}_j- \mb{w}_j)= \mb{0}_V$. Since each $\mb{v}_j- \mb{w}_j \in W_j$, we have $\mb{v}_j= \mb{w}_j$.
2 $\to$ 3: Clearly $\Span( \cup_{j=1}^k \gamma_j)=V$. To see $\cup_{j=1}^k \gamma_j$ is LID, say they are $\mb{w}_1, \ldots \mb{w}_n$, the 1st $m$ are in $\gamma_1$, and $\sum_{s=1}^na_s \mb{w}_s= \mb{0}_V$.
Then $\sum_{s=1}^ma_s \mb{w}_s \in W_1$ and sums of groups of the rest are in $W_2, \ldots W_k$, resp. Thus each group sum is $\mb{0}_V$, in particular $\sum_{s=1}^ma_s \mb{w}_s= \mb{0}_V$. But since $\gamma_1$ is LID, we have $a_j=0$ for $j= \ot{m}$. The other coefficients are likewise 0.
3 $\to$ 4 is immediate. 4 $\to$ direct sum: clearly $V= \sum_{j=1}^kW_j$. If $\mb{v}\in W_1 \cap \sum_{j=2}^kW_j$, then it is a linear combination, both of $\gamma_1$ and $\cup_{j=2}^k \gamma_j$. Since $\cup_{j=1}^k \gamma_j$ is a basis, this is not possible unless all coefficients are 0, ie $\mb{v}_j= \mb{0}_V$.
\end{proof}
\begin{prop}\label{p:sum-espace}
A linear operator $T$ on a finite dimensional vector space $V$ is diagonalizable if and only if $V$ is the direct sum of its eigenspaces.
\end{prop}
\begin{proof}%easy, short
By Proposition \ref{p:diagze-mult}, $T$ is diagonalizable if and only if $\cup_{j=1}^k \beta_j$ is an eigenvector basis for $V$, where each $\beta_j$ is a basis for an eigenspace. By Lemma \ref{l:directsum}, this is if and only if $V$ is the direct sum of those eigenspaces.
\end{proof}
\begin{exer}\label{e:split-ut}
Let $A \in F^{n \times n}$. Analogous results hold for linear operators.
\begin{itemize}
\item For any eigenvalue $\lambda$ of $A$, that eigenspace is the same as the eigenspace of $1/ \lambda$ for $A^{-1}$. $A$ is thus diagonalizable if and only if $A^{-1}$ is.
\item $A$ and $A^\top$ have the same set of eigenvalues (with the same multiplicities), and their eigenspaces for the same eigenvalue have the same dimension. $A$ is thus diagonalizable if and only if $A^\top$ is.
\item $A$ is similar to an upper triangular matrix if and only if its characteristic polynomial splits.
\item If the characteristic polynomial of $A$ splits, and it has (not necessarily distinct) eigenvalues $\lambda_1, \ldots, \lambda_n$, then $\tr A= \sum_{j=1}^n \lambda_j$ and $\det A= \prod_{j=1}^n \lambda_j$.
\end{itemize}
\end{exer}
\begin{proof}
Part 1: We can use the matrix version of Exercise \ref{e:eig-poly}, but it is also easy to see directly that $\{ \mb{v}\in F^n: A \mb{v}= \lambda \mb{v}\} = \{ \mb{v}\in F^n: A^{-1}\mb{v}= \mb{v}/ \lambda \}$, and that any basis of $F^n$ of eigenvectors of $A$ also consists of eigenvectors of $A^{-1}$.\\
Part 2: $A$ and $A^\top$ have the same characteristic polynomial (Lemma \ref{l:chrpoly-degn}), so they have the same eigenvalues with the same multiplicities, and 1 splits if and only if the other does.
Using the dimension formula (Lemma \ref{l:dim-formula}) and Corollary \ref{c:rank-transpose}, we have $\dim \ker(A^\top - \lambda I_n)= \dim \ker(A- \lambda I_n)^\top =n- \rank(A- \lambda I_n)^\top =n- \rank(A- \lambda I_n)= \dim \ker (A- \lambda I_n)$.
 $A$ is thus diagonalizable if and only if $A^\top$ is by Proposition \ref{p:diagze-mult}.\\
Part 3: "$\to$" If $A$ is similar to upper triangular $U$, then they have the same characteristic polynomial (Lemma \ref{l:chr-poly}), which is $t \mapsto \prod_{j=1}^n(u_{jj}-t)$, which splits.
"$\leftarrow$" We use induction on $n$. The base case $n=1$ is immediate. When $n>1$, the fact that the characteristic polynomial splits means we can pick an eigenvalue $\lambda$ and eigenvector $\mb{v}_1$ of $A$.
Extend it to a basis $mb{v}_1, \ldots, mb{v}_n$ of $F^n$ (Corollary \ref{c:extend-LID2basis}). If $P$ is a matrix with these vectors as columns, then
\[ P^{-1}AP= \tbtm{\lambda}{\mb{0}_{n-1}}{*}{B}\]
The characteristic polynomial of $A$ is the characteristic polynomial of the RHS (Lemma \ref{l:chr-poly}), which is $( \lambda -t)$ times the characteristic polynomial of $B$ (Exercise \ref{e:block-det}), which must split.
By inductive hypothesis, there exists $R \in F^{(n-1) \times(n-1)}$ such that $R^{-1}BR$ is upper triangular, and
\[ \tbtm{1}{\mb{0}_{n-1}}{\mb{0}_{n-1}^\top}{R}^{-1}P^{-1}AP \tbtm{1}{\mb{0}_{n-1}}{\mb{0}_{n-1}^\top}{R}= \tbtm{\lambda}{\mb{0}_{n-1}}{*}{R^{-1}BR}\]
So $A$ is similar to an upper triangular matrix.
Part 4: By part 3, there exists $P$ invertible such that $U:= P^{-1}AP$ is upper triangular, with $\lambda_1, \ldots, \lambda_n$ along its main diagonal. By Lemma \ref{l:chr-poly}, $\tr A= \tr U= \sum_{j=1}^n \lambda_j$, and $\det A= \det U= \prod_{j=1}^n \lambda_j$.
\end{proof}
\begin{defn}
Let $V$ be a vector space, $W_1,W_2$ subspaces such that $V=W_1 \oplus W_2$. The \emph{projection} of $V$ on $W_1$ along $W_2$ is $T:V \to V$, $\mb{x}\mapsto \mb{x}_1$, where $\mb{x}= \mb{x}_1+ \mb{x}_2$, $\mb{x}_1 \in W_1$, $\mb{x}_2 \in W_2$.
\end{defn}
\begin{exer}\label{e:projection}
In the setting of the above definition
\begin{itemize}
\item $T$ is linear, $W_1= \{ \mb{x}\in V:T( \mb{x})= \mb{x}\} = \img T$, $W_2= \ker T$. In particular, $V= \img T \oplus \ker T$, and no ambiguity arises from simply calling $T$ a projection.\footnote{It is the projection of $V$ on its image along its kernel}
\item $T \in \cl{L}(V)$ is a projection if and only if $T^2=T$.
\end{itemize}
\end{exer}
\begin{proof}%rating 1/2
Part 1: Fix $\mb{x}, \mb{y}\in V$, $c \in F$. Write $\mb{x}= \mb{x}_1+ \mb{x}_2$, $\mb{y}= \mb{y}_1+ \mb{y}_2$, where $\mb{x}_1, \mb{y}_1 \in W_1$, $\mb{x}_2, \mb{y}_2 \in W_2$.
Then $c \mb{x}+ \mb{y}=(c \mb{x}_1+ \mb{y}_1)+(c \mb{x}_2+ \mb{y}_2)$, where $c \mb{x}_1+ \mb{y}_1 \in W_1$, $c \mb{x}_2+ \mb{y}_2 \in W_2$.
Thus, $T(c \mb{x}+ \mb{y})=c \mb{x}_1+ \mb{y}_1= c \cdot T( \mb{x})+T( \mb{y})$, so $T$ is linear.
Next, the 3 inclusions $W_1 \subseteq \{ \mb{x}\in V:T( \mb{x})= \mb{x}\} \subseteq \img T \subseteq W_1$ are all immediate, as is $W_2 \subseteq \ker T$.
If $\mb{x}\in \ker T$, write $\mb{x}= \mb{x}_1+ \mb{x}_2$, $\mb{x}_1 \in W_1$, $\mb{x}_2 \in W_2$. Then $\mb{0}_V=T( \mb{x})= \mb{x}_1$, so $\mb{x}= \mb{x}_2 \in W_2$.
Part 2: "$\to$" For $\mb{x}\in V$, write $\mb{x}= \mb{x}_1+ \mb{x}_2$, $\mb{x}_1 \in W_1$, $\mb{x}_2 \in W_2$. Then $T^2( \mb{x})=T( \mb{x}_1)= \mb{x}_1=T( \mb{x})$.
"$\leftarrow$" For $\mb{x}\in V$, write $\mb{x}=T( \mb{x})+ \{ \mb{x}-T( \mb{x})\}$. Since $T \in \cl{L}(V)$ and $T^2=T$, $T \big( \mb{x}-T( \mb{x}) \big)= \mb{0}_V$, so $V= \img T+ \ker T$.
If $\mb{x}\in \img T \cap \ker T$, then $\mb{x}=T( \mb{y})$ for some $\mb{y}\in V$, and $\mb{0}_V=T( \mb{x})=T^2( \mb{y})=T( \mb{y})= \mb{x}$. Thus, $\img T \cap \ker T= \{ \mb{0}_V \}$, and in fact $V= \img T \oplus \ker T$.
\end{proof}

\section{Matrix limits and Markov chains*}
We study $\lim_{m \to \infty}A^m$ for $A \in \bb{C}^{n \times n}$, which has applications in the natural and social sciences.
Recall that a sequence $(z_n)$ in $\bb{C}$ converges if $( \Re z_n)$ and $( \Im z_n)$ do. A sequence $(A_m)$ in $\bb{C}^{n \times p}$ converges if each entry does.
Using $\lim_{n \to \infty}(ca_n+b_n)=c \lim_{n \to \infty}a_n+ \lim_{n \to \infty}b_n$ for sequences $(a_n)$ and $(b_n)$ in $\bb{C}$ when the limits exist, we have:
\begin{lem}\label{l:lin-lmt}
For a sequence $(A_m)$ in $\bb{C}^{n \times p}$, and compatible matrices $B,C$, we have $\lim_{m \to \infty}BA_m=B \lim_{m \to \infty}A_m$ and $\lim_{m \to \infty}A_mC=( \lim_{m \to \infty}A_m)C$. In particular, when $A,Q \in \bb{C}^{n \times n}$ and $Q$ is invertible, then $\lim_{m \to \infty}(QAQ^{-1})^m= Q( \lim_{m \to \infty}A^m)Q^{-1}$.
\end{lem}
This is a set we frequently encounter:
\begin{lem}\label{l:mat-power0}
Put $S:= \{ z \in \bb{C}: \lim_{n \to \infty}z^n \T{ exists}\} = B_1(0) \cup {1}$. For $A \in \bb{C}^{n \times n}$, if both:
\begin{itemize}
\item Each eigenvalue of $A$ is in $S$.
\item $A$ is diagonalizable.
\end{itemize}
Then $\lim_{m \to \infty}A^m$ exists. Condition 1 is also necessary.
\end{lem}
We later show that if we replace condition 2 with: $\dim E_1$ is the multiplicity of eigenvalue $E_1$, then the 2 conditions are also necessary. As an illustration,
\[ \tbtm{1}{0}{1}{1}^m= \tbtm{1}{0}{m}{1}\]
\begin{proof}
If $|z|<1$, then $\lim_{n \to \infty}z^n=0$. If $|z|>1$, then $\lim_{n \to \infty}z^n= \infty$. If $|z|=1$, write $z=e^{si}$ for some $0 \leq s <2 \pi$.
Then $\lim_{n \to \infty}z^n= \lim_{n \to \infty}e^{nsi}$, which only exists when $s=0$ since that is the only value of $s$ where $(ns \bmod 2 \pi)$ converges.
Use Lemma \ref{l:diagze-basis} to write $D=Q^{-1}AQ$ $\leftrightarrow$ $A=QDQ^{-1}$ where $Q$ has eigenvectors of $A$ as its columns and $D$ is diagonal with the corresponding eigenvalues $\lambda_j$ along its main diagonal.
By Lemma \ref{l:lin-lmt}, $\lim_{m \to \infty}(QDQ^{-1})^m= Q( \lim_{m \to \infty}D^m)Q^{-1}$, which exists if and only if each $\lambda_j \in S$.
Condition 1 is necessary since if $A \mb{v}= \lambda \mb{v}$ for some $\lambda \notin S$ and nonzero $\mb{v}\in \bb{C}^n$, then $( \lim_{m \to \infty}A^m) \mb{v}= \lim_{m \to \infty}A^m \mb{v}= \lim_{m \to \infty}\lambda^m \mb{v}$, which does not exist, so neither does $\lim_{m \to \infty}A^m$.
\end{proof}
\begin{defn}
$P \in \bb{R}^{n \times n}$ is a \emph{transition matrix} (or stochastic matrix) If each entry $p_{jk}\geq 0$ and $P \mb{1}_n= \mb{1}_n$ (all rows sum to 1).
Each $p_{jk}$ is interpreted as the probability to move from state j to state k in 1 stage.\footnote{Friedberg uses $P^\top$ as the transition matrix, hence his column sums are 1 and he interprets $p_{jk}$ as state k to j}
$\mb{u}\in \bb{R}^n$ is a \emph{probability vector} if each $u_j \geq 0$ and $\mb{1}_n^\top \mb{u}=1$.\footnote{the set of all probability vectors is often denoted $\Delta^{n-1}$, the n-1 dimensional probability simplex}
\end{defn}
\begin{lem}\label{l:tm-def}
Let $P,Q \in \bb{R}^{n \times n}$ be transition matrices, $\mb{u}\in \bb{R}^n$ a probability vector.
\begin{itemize}
\item $PQ$ and $P^m$ for $m \in \bb{Z}_{\geq 0}$ are transition matrices.
\item If $u_j$ is interpreted as the initial probability of being in state j, and $P$ is interpreted as in the definition, then the jth entry of $\mb{u}^\top P^m$ is the probability of being in state j after m stages.
\item If $L:= \lim_{m \to \infty}P^m$, then $\mb{u}^\top LP= \mb{u}^\top L$. The jth entry of $\mb{u}^\top L$ is the long term probability of being in state j.
\end{itemize}
\end{lem}
\begin{proof}%east, short
Part 1: Clearly each $(PQ)_{jk}\geq 0$ and $PQ \mb{1}_n= P \mb{1}_n= \mb{1}_n$. It follows that each $P^m$ is a transition matrix by induction.
Part 2: For the probability of being in state j after 1 stage, consider initially being in state k for $k= \ot{n}$. In symbols, this is $\sum_{k=1}^nu_kp_{jk}$, the jth entry of $\mb{u}^\top P$. Applying this m times, we get $\mb{u}^\top P^m$.
Part 3: Applying Lemma \ref{l:lin-lmt} twice, we have $\mb{u}^\top( \lim_{m \to \infty}P^m)P= \mb{u}^\top \lim_{m \to \infty}P^{m+1}= \mb{u}^\top L$.
\end{proof}
\begin{defn}
A \emph{stochastic process} is a sequence $(X_m \in \cl{S}:m \in \bb{Z}_{\geq 0})$ of random variables. $\cl{S}$ is the set of possible states.
A \emph{Markov process} is a stochastic process where the transition probability to 1 state depends only on the previous state and not on the stage.
A \emph{Markov chain} is a Markov process where $\cl{S}$ is finite.\footnote{more often, countable}
\end{defn}
For later: using refined version of Lemma \ref{l:mat-power0}, if $P$ is a transition matrix, then $\lim_{m \to \infty}P^m$ exists if and only if all its eigenvalues are in $S$. An illustrative example is
\[ \lim_{m \to \infty}\tbtm{0}{1}{1}{0}^m \]
does not exist (the sequence alternates between that matrix and $I 2$) since -1 is an eigenvalue.
Even when $\lim_{m \to \infty}P^m$ exists, it can be hard to compute. We study conditions where it is easily computed.
\begin{defn}
A transition matrix $P$ is \emph{regular} if there exists $m \in \bb{Z}^+$ such that each entry of $P^m$ is positive.\footnote{A Markov chain with a regular transition matrix is irreducible and aperiodic}
\end{defn}
\begin{example}
A state j of a Markov chain is \emph{absorbing} if it can never be left once entered. The jth row of the transition matrix (and its powers) is then $\mb{e}_j$, so it is not regular.
The entire chain is called \emph{absorbing} if it is possible to go from any state into an absorbing one. We do not study absorbing states or chains.
\end{example}
\begin{defn}
For $A \in \bb{C}^{n \times p}$, its \emph{row sum} is $\rho(A):= \max_{j= \ot{n}}\rho_j(A)$, where $\rho_j(A):= \sum_{k=1}^p|a_{jk}|$. Likewise, its \emph{column sum} is $\nu(A):= \max_{k= \ot{p}}\nu_j(A)$, where $\nu_j(A):= \sum_{j=1}^n|a_{jk}|$.\footnote{$\rho(A)= \max \{ \| A \mb{x}\|_\infty : \| \mb{x}\|_\infty \leq 1 \}$ and $\nu(A)= \max \{ \| A \mb{x}\|_1 : \| \mb{x}\|_1 \leq 1 \}$}
\end{defn}
\begin{thm}[Gershgorin Circle Theorem]
For $A \in \bb{C}^{n \times n}$, its jth \emph{Gershgorin disc} is $C_j:= \{ z \in \bb{C}: |z-a_{jj}| \leq \rho_j(A)-|a_{jj}| \}$. All eigenvalues of $A$ are contained in $\cup_{j=1}^nC_j$.
\end{thm}
While not ours, 1 use case for this theorem is to show if 0 is not in any Gershgorin disc, then $A$ is invertible (Exercise \ref{e:eig-poly}).
\begin{proof}%easy, short
Let $\lambda$ be an eigenvalue, pick nonzero $\mb{v}\in \bb{C}^n$ such that $A \mb{v}= \lambda \mb{v}$. Pick j that maximizes $|v_j|$, then $v_j \neq 0$, and $\sum_{k=1}^na_{jk}v_k= \lambda v_j$.
Then $\lambda -a_{jj}= \sum_{k \neq j}a_{jk}v_k/v_j$, so $| \lambda -a_{jj}| \leq \sum_{k \neq j}|a_{jk}|= \rho_j(A)- |a_{jj}|$.
\end{proof}
\begin{corollary}\label{c:gershgorin}
For $A \in \bb{C}^{n \times n}$, all its eigenvalues $\lambda$ satisfy $| \lambda| \leq \rho(A) \wedge \nu(A)$.
For a transition matrix $P \in \bb{R}^{n \times n}$, 1 is an eigenvalue, and all eigenvalues $\lambda$ satisfy $| \lambda| \leq 1$.
\end{corollary}
\begin{proof}%easy, short
By Gershgorin's circle theorem, there exists $j \in \{ \ot{n}\}$ such that $| \lambda| \leq | \lambda -a_{jj}|+|a_{jj}| \leq \rho_j(A)-|a_{jj}|+|a_{jj}| \leq \rho(A)$.
By Exercise \ref{e:split-ut}, $A$ and $A^\top$ have the same set of eigenvalues, so we also have $| \lambda| \leq \rho(A^\top)= \nu(A)$.
The assertions about transition matrices follow immediately from $P \mb{1}_n= \mb{1}_n$ and $\rho(P)=1$.
\end{proof}
\begin{prop}\label{p:rowsum-evec}
If $A \in \bb{R}^{n \times n}$ with each $a_{jk}>0$ and $\lambda$ is an eigenvslue with $| \lambda|= \rho(A)$, then $\lambda= \rho(A)$ and the eigenspace is $E_{\rho(A)}= \Span \{ \mb{1}_n \}$.
If $\lambda$ is an eigenvslue with $| \lambda|= \nu(A)$, then $\lambda= \nu(A)$ and $\dim E_{\nu(A)}=1$.
\end{prop}
\begin{proof}
We start the same way as we did in proving Gershgorin's Circle Theorem. Let $\mb{v}$ be an eigenvector for $\lambda$, then $A \mb{v}= \lambda \mb{v}$.
Pick j to maximize $|v_j|$. Then $v_j \neq 0$, and $\sum_{k=1}^na_{jk}v_k= \lambda v_j$. Then $| \lambda| \cdot|v_j|=| \lambda v_j|=| \sum_{k=1}^na_{jk}v_k| \leq \sum_{k=1}^na_{jk}|v_k| \leq |v_j| \sum_{k=1}^na_{jk}\leq |v_j| \rho(A)$.
But by hypothesis $| \lambda|= \rho(A)$, and all the inequalities in this chain are actually equalities. In particular, $| \sum_{k=1}^na_{jk}v_k|= \sum_{k=1}^na_{jk}|v_k|$ and each $a_{jk}>0$ means all the $\{ v_k \}$ point in the same direction in the complex plane, so each $v_k=|v_k|e^{si}$ for some universal $0 \leq s<2 \pi$.
Next, $\sum_{k=1}^na_{jk}|v_k|=|v_j| \sum_{k=1}^na_{jk}$ means $|v_1|= \ldots=|v_n|$ and so $v_1= \ldots=v_n$. This shows that $\mb{1}_n$ is an eigenvector for $\lambda$, and (since $\mb{v}\in E_\lambda$ was arbitrary) that $E_\lambda= \Span \{ \mb{1}_n \}$.
The final inequality that is actually an equality is $\sum_{k=1}^na_{jk}= \rho(A)$. Combine this with $\sum_{k=1}^na_{jk}v_k= \lambda v_j$ and $v_1= \ldots=v_n$ to get $\lambda= \rho(A)$ and thus $\mb{v}\in \bb{R}^n$.
We now apply this result to $A^\top$ (Exercise \ref{e:split-ut}) to get the statement with $\nu(A)$.
\end{proof}
\begin{corollary}\label{c:rowsum-evec}
If $P$ is a transition matrix and each $p_{jk}>0$, then 1 is an eigenvalue with eigenspace $E_1= \Span \{ \mb{1}_n \}$, and all other eigenvalues $\lambda$ satisfy $| \lambda|<1$. The same conclusions hold if $P$ is regular.
\end{corollary}
\begin{proof}%easy, short
When each $p_{jk}>0$, all the conclusions follow from Proposition \ref{p:rowsum-evec} since by definition $\rho(P)=1$ for transition matrices.
If $P$ is regular, pick $m \in \bb{Z}^+$ such that all entries of $P^m$ are positive.
Then so are all entries of $P^{m+1}$ ($P^{m+1}=PP^m$, and all rows of $P$ sum to 1), and both are transition matrices (Lemma \ref{l:tm-def}), so Proposition \ref{p:rowsum-evec} applies to both.
In particular, if $\lambda$ is an eigenvalue of $P$ with $| \lambda|=1$ (rule out $| \lambda|>1$ with Corollary \ref{c:gershgorin}) and eigenvector $\mb{v}$, then $\lambda^m$ and $\lambda^{m+1}$ are eigenvalues of $P^m$ and $P^{m+1}$ with eigenvector $\mb{v}$.
Since $| \lambda^m|=| \lambda^{m+1}|=1$ we actually have $\lambda^m= \lambda^{m+1}=1$ and hence $\lambda=1$. Further, $\mb{v}\in \Span \{ \mb{1}_n \}$.
\end{proof}
\begin{prop}
If $P$ is a diagonalizable regular transition matrix, then $L:= \lim_{m \to \infty}P^m$ exists.\footnote{diagonalizability is not needed. All other assertions remain true with their proofs unchanged} Further:
\begin{itemize}
\item $L$ is a transition matrix, and $PL=LP=L$.
\item All rows of $L$ are identical, and the eigenspace of $P^\top$ for eigenvalue 1 is $\Span \{ L^\top \mb{e}_1 \}$.
\item For all probability vectors $\mb{u}$, we have $\lim_{m \to \infty}\mb{u}^\top P^m= \mb{e}_1^\top L$. We call $L^\top \mb{e}_1$ the \emph{stationary vector} (or fixed probability vector) of $P$. 
\end{itemize}
\end{prop}
Based on part 2, a good way to find the stationary vector is to sole $(P^\top -I_n) \mb{v}= \mb{0}_n$ for $\mb{v}$.
\begin{proof}%easy, short
The limit exists by Corollary \ref{c:rowsum-evec} and Lemma \ref{l:mat-power0}.\\
Part 1: Each $P^m$ is a transition matrix (Lemma \ref{l:tm-def}), so each $L_{jk}= \lim_{m \to \infty}(P^m)_{jk}\geq 0$ and $\sum_{k=1}^nL_{jk}= \sum_{k=1}^n \lim_{m \to \infty}(P^m)_{jk}= \lim_{m \to \infty}\sum_{k=1}^n(P^m)_{jk}=1$.
Also, $PL=LP= \lim_{m \to \infty}P^{m+1}=L$ by Lemma \ref{l:lin-lmt}.
Part 2: Since $PL=L$, all columns of $L$ are in the eigenspace of $P$ with eigenvalue 1. They are thus each a scalar multiple of $\mb{1}_n$, so each row is identical.
$LP=L$ $\leftrightarrow$ $P^\top L^\top =L^\top$ (Exercise \ref{e:transpose-inv}), so al rows of $L$ (WLOG the 1st) are in the eigenspace of $P^\top$ with eigenvalue 1, so $L^\top \mb{e}_1$ spans that space.
Part 3: By Lemma \ref{l:lin-lmt}, $\lim_{m \to \infty}\mb{u}^\top P^m= \mb{u}^\top L= \mb{e}_1^\top L$ since $\mb{u}$ is a probability vector and all rows of $L$ are the same.
\end{proof}
\begin{example}[Hardy Weinberg Law]
Let the initial distribution of genotypes GG, Gg, and gg in a population be given by a probability vector $\mb{v}\in \bb{R}^3$. Assume the population is big, and that gender/ mating/ survival is random with respect to genotype.
Then in just 1 generation, the distribution will reach equilibrium with probability vector $\mb{u}:= \big( p^2,2p(1-p),(1-p)^2 \big)^\top$, where $p:=v_1+v_2/2$ is the frequency of allele G ($1-p=v_2/2+v_3$ is the frequency of allele g).
The initial transition matrix based on $\mb{v}$ is
\[ P_0:= \begin{pmatrix} v_1+v_2/2 & v_2/2+v_3 & 0 \\
v_1+v_2/4 & (v_1+v_2+v_3)/2 & v_2/4+v_3/2 \\
0 & v_1+v_2/2 & v_2/2+v_3 \end{pmatrix}= \begin{pmatrix} p & 1-p & 0 \\
p/2 & 1/2 & (1-p)/2 \\
0 & p & 1-p \end{pmatrix}\]
Each entry is computed by $P_0$(offspring has x|1 parent has y)=$\sum_zP_0(z)P_0$(offspring has x|parents have y,z) for $x,y,z \in \{ GG, Gg, gg \}$.
In particular, the matrix $P_0$ depends on $\mb{v}$ only thru $p$. Provided $p$ does not change per generation, the transition matrix from any generation to the next is $P=P_0$.
The distribution for generation 1 is $\mb{v}^\top P_0= \big( p(v_1+v_2/2),(1-p)v_1+v_2/2+pv_3,(1-p)(v_2/2+v_3) \big) = \big( p^2,2p(1-p),(1-p)^2 \big) = \mb{u}^\top$.
Since $u_1+u_2/2=p$, we have $P_1=P_0$ and $\mb{u}^\top P_1= \mb{u}^\top$.
\end{example}

\section{Invariant subspaces}
We have seen that if $T$ has eigenvector $\mb{v}$, then $T$ maps $\Span \{ \mb{v}\}$ into itself. This kind of behavior is important for linear operators.
\begin{defn}
Let $V$ be a vector space and $T \in \cl{L}(V)$. A subspace $W$ is T-\emph{invariant} if $T(W) \subseteq W$.
If $\mb{x}\in V$, $\mb{x}\neq \mb{0}_V$, call $\Span \{ T^m( \mb{x}): m \in \bb{Z}_{\geq 0}\}$ the T-\emph{cyclic} subspace generated by $\mb{x}$.
If $W$ is a subspace of $V$, define the \emph{restriction} of $T$ to $W$ by $T_W:W \to V$, $\mb{v}\mapsto T( \mb{v})$.
\end{defn}
\begin{exer}\label{e:smallest-tcyclic}
These subspaces are always T-invariant: $\{ \mb{0}_V \}$, $V$, $\img T$, $\ker T$, any eigenspace $E_\lambda$ and the T-cyclic subspace generated by $\mb{x}$. It is the smallest T-invariant subspace containing $\mb{x}$.
We have $T_W \in \cl{L}(V,W)$, and when $W$ is T-invariant, $T_W \in \cl{L}(W)$.
\end{exer}
\begin{proof}
Let $W_x$ denote this T-cyclic subspace. For any $m \in \bb{Z}_{\geq 0}$ we have $T \circ T^m( \mb{x})= T^{m+1}( \mb{x}) \in W_x$, that is $T(W_x) \subseteq W_x$, so it is T-invariant.
If $W$ is any T-invariant subspace containing $\mb{x}$ then it must contain each $T^m( \mb{x})$ by induction, so it contains their span $W_x$ (Lemma \ref{l:span-minimal}).
\end{proof}
Naturally, $T_W$ inherits many properties of $T$.
\begin{lem}\label{l:restrict-chrpoly}
If $V$ is a finite dimensional vector space, $T \in \cl{L}(V)$, and $W$ is a T-invariant subspace, then the characteristic polynomial of $T_W$ divides that of $T$.
\end{lem}
\begin{proof}
Let $\beta:= \{ \mb{v}_1, \ldots, \mb{v}_n \}$ be a basis for $V$ where the 1st m vectors are a basis $\gamma$ for $W$ (Corollary \ref{c:extend-LID2basis}). For $j= \ot{m}$, the coordinates of $T( \mb{v}_j)=T_W( \mb{v}_j)$ with respect to $\beta$ are 0 except for possibly the 1st m, so
\[ [T]_\beta = \tbtm{[T_W]_\gamma}{O}{*}{B}\]
so $\det ([T]_\beta -tI_n)= \det ([T_W]_\gamma -tI_m) \det (B-tI_{n-m})$ (Exercise \ref{e:block-det}).
\end{proof}
The characteristic polynomial of $T_W$ thus gives insight into that of $T$. If $W$ is also T-cyclic, then this polynomial is easy to compute.
\begin{prop}\label{p:tcyclic-poly}
If $V$ is a finite dimensional vector space, $T \in \cl{L}(V)$, $\mb{v}\in V$, and $W$ is the T-cyclic subspace generated by $\mb{v}$. Pick the smallest $m \in \bb{Z}^+$ such that $\beta_m:= \{ T^j( \mb{v}):j=0, \ldots,m \}$ is LD.
Then $\beta_{m-1}$ is a basis for $W$, and if $\sum_{j=0}^ma_jT^j( \mb{v})= \mb{0}_V$ with $a_m=1$, then $t \mapsto (-1)^m \sum_{j=0}^ma_jt^j$ is the characteristic polynomial of $T_W$.
\end{prop}
\begin{proof}%rating 1/2
Such an $m$ exists and $m \leq \dim V$ (Corollary \ref{c:bases-n-same}). Put $Z:= \Span \beta_{m-1}$. Then $\beta_{m-1}$ is a basis for $Z$ and $T^m( \mb{v}) \in Z$ (Lemma \ref{l:span-LD}).
We claim that $Z$ is T-invariant. Indeed, if $\mb{w}\in Z$, then $\mb{w}= \sum_{j=0}^{m-1}b_jT^j( \mb{v})$ for some $b_0, \ldots, b_{m-1}\in F$, so $T( \mb{w})= \sum_{j=0}^{m-1}b_jT^{j+1}( \mb{v})$.
Since $\mb{v}\in Z$ too, $W \subseteq Z$ (Exercise \ref{e:smallest-tcyclic}), so in fact $W=Z$. Put $\beta:= \beta_{m-1}$. Then
\[ [T_W]_\beta= \tbtm{\mb{0}_{m-1}^\top}{I_{m-1}}{-a_0}{\begin{matrix}\vdots \\ -a_{m-1}\end{matrix}}\]
The characteristic polynomial has the stated form by induction on m. For the base case $m=2$, $\det ([T_W]_\beta -tI_2)=-t(-a_1-t)+a_0=(-1)^2 \sum_{j=0}^2a_jt^j$ ($a_2=1$).
For $m>2$, cofactor expansion along the 1st row gives $\det ([T_W]_\beta -tI_m)=-t \det \{ ( \overline{[T_W]_\beta})_{11}-tI_{m-1}\})+(-1)^{m+1}(-a_0)=-t(-1)^{m-1} \sum_{j=1}^{m-1}a_jt^{j-1}+(-1)^ma_0= (-1)^m \sum_{j=0}^ma_jt^j$, where we used the inductive hypothesis on the 1,1 cofactor.
\end{proof}
This is the key tool that proves the famous result that a linear operator/ matrix is annihilated by its own characteristic polynomial.
\begin{thm}[Cayley Hamilton]
If $V$ is a finite dimensional vector space, $T \in \cl{L}(V)$, and $f$ is its characteristic polynomial, then $f(T) \equiv \mb{0}_V$.
\end{thm}
\begin{proof}%easy, short
Fix $\mb{v}\in V$. If $\mb{v}= \mb{0}_V$, then $f(T)( \mb{v})= \mb{0}_V$. Else, let $W$ be the T-cyclic subspace generated by $\mb{v}$. Construct the characteristic polynomial $g(t):=(-1)^m \sum_{j=0}^ma_jt^j$ of $T_W$ as in Proposition \ref{p:tcyclic-poly}, which also states that $g(T)( \mb{v})=g(T_W)( \mb{v})= \mb{0}_V$.
By Lemma \ref{l:restrict-chrpoly}, $g$ divides $f$, so $f(T)( \mb{v})= \mb{0}_V$.
\end{proof}
If we can decompose a finite dimensional vector space into a direct sum of T-invariant subspaces, then we can infer the behavior of $T$ based on those direct summands.
\begin{defn}
The \emph{direct sum} of square matrices $B_1, \ldots,B_k$ is the block diagonal matrix they form, denoted $\bigoplus_{j=1}^kB_k$.
\end{defn}
\begin{exer}\label{e:dirsum-mat}
If $V$ is a vector space with $\dim V=n$ then $T \in \cl{L}(V)$ is diagonalizable if and only if $V$ can be written as the direct sum of 1 dimensional T-invariant subspaces.\\
If $W_1, \ldots,W_k$ are T-invariant subspaces with ordered bases $\beta_1, \ldots, \beta_k$, $\beta:= \cup_{j=1}^k \beta_j$, $V= \bigoplus_{j=1}^kW_j$, then $[T]_\beta= \bigoplus_{j=1}^k[T_{W_j}]_{\beta_j}$.
\end{exer}
\begin{proof} %easy, short
Part 1: "$\to$" Let $\{ \mb{v}_1, \ldots, \mb{v}_n \}$ be a basis of $V$ of eigenvectors of $T$. For $j= \ot{n}$, put $W_j:= \Span \{ \mb{v}_j \}$. Then $V= \bigoplus_{j=1}^nW_j$, each $W_j$ is T-invariant, and $\dim W_j=1$.
"$\leftarrow$" If $V= \bigoplus_{j=1}^nW_j$, each $W_j$ is T-invariant, and $\dim W_j=1$, then for $j= \ot{n}$, pick nonzero $\mb{v}_j \in W_j$. Then $T( \mb{v}_j) \in W_j$, so $T( \mb{v}_j)= \lambda_j \mb{v}_j$ for some $\lambda_j \in F$.\\
Part 2: $\beta$ is indeed a basis (Lemma \ref{l:directsum}).%We show it for $k=2$, the general case follows by induction (since $\bigoplus_{j=1}^{k-1}W_j$ is T-invariant and $V=( \bigoplus_{j=1}^{k-1}W_j) \oplus W_k$). 
For $j= \ot{k}$ and any $\mb{v}\in W_j$, we have $T( \mb{v}) \in W_j$, so it is a linear combination of $\beta_j$. It follows that $[T]_\beta= \bigoplus_{j=1}^k[T_{W_j}]_{\beta_j}$.
\end{proof}
\begin{lem}
If $V$ is a vector space with $\dim V=n$, $T \in \cl{L}(V)$, $V= \bigoplus_{j=1}^kW_j$ where each $W_j$ is T-invariant, and $f,f_j$ are the characteristic polynomials of $T,T_{W_j}$, then $f= \prod_{j=1}^kf_j$. In particular, $\det T= \prod_{j=1}^k \det T_{W_j}$.
\end{lem}
This is close to Lemma \ref{l:restrict-chrpoly}. Unfortunately, given a T-invariant subspace $W_1$, at this point we do not 'know' there exists another $W_2$ such that $V=W_1 \oplus W_2$.
\begin{proof}%easy, short
Set things up as in Exercise \ref{e:dirsum-mat}. Then $[T]_\beta= \bigoplus_{j=1}^k[T_{W_j}]_{\beta_j}$, so $\det([T]_\beta -tI_n)= \det \big\{ \bigoplus_{j=1}^k([T_{W_j}]_{\beta_j}-tI_{m_j} \big\} = \prod_{j=1}^k \det ([T_{W_j}]_{\beta_j}-tI_{m_j})$ by Exercise \ref{e:block-det} and induction, where $m_k$ is the number of vectors in $\beta_k$.
Take $t=0$ to get the claim about determinants.
\end{proof}
%\{}